\documentclass[twocolumn,11pt,a4paper]{article}

\usepackage[utf8]{inputenc}
\usepackage[T1]{fontenc}
\usepackage{lmodern}
\usepackage[margin=1in,columnsep=0.6cm]{geometry}
\usepackage{amsmath,amssymb,amsthm,mathtools}
\usepackage{thmtools}
\usepackage{enumitem}
\usepackage{microtype}
\usepackage{booktabs}
\usepackage{array}
\usepackage{graphicx}
\usepackage{hyperref}
\usepackage{cleveref}
\usepackage{xcolor}
\usepackage[numbers,sort&compress]{natbib}
\usepackage{tcolorbox}

\hypersetup{colorlinks=true,linkcolor=blue!50!black,
            citecolor=green!50!black,urlcolor=blue!60!black}

%---- theorem environments ----
\declaretheoremstyle[spaceabove=5pt,spacebelow=5pt,
  headfont=\bfseries,bodyfont=\itshape,postheadspace=1em]{thmstyle}
\declaretheoremstyle[spaceabove=5pt,spacebelow=5pt,
  headfont=\bfseries,bodyfont=\normalfont,postheadspace=1em]{defstyle}

\declaretheorem[style=thmstyle,name=Theorem,numberwithin=section]{theorem}
\declaretheorem[style=thmstyle,name=Lemma,sibling=theorem]{lemma}
\declaretheorem[style=thmstyle,name=Proposition,sibling=theorem]{proposition}
\declaretheorem[style=thmstyle,name=Corollary,sibling=theorem]{corollary}
\declaretheorem[style=defstyle,name=Definition,sibling=theorem]{definition}
\declaretheorem[style=defstyle,name=Axiom,sibling=theorem]{axiom}
\declaretheorem[style=defstyle,name=Remark,sibling=theorem]{remark}
\declaretheorem[style=defstyle,name=Example,sibling=theorem]{example}

%---- notation ----
\newcommand{\kB}{k_{\mathrm{B}}}
\newcommand{\Sfloor}{S_{\flat}}
\newcommand{\Gcond}{\mathcal{G}}
\newcommand{\HolL}{H_\ell}
\newcommand{\drugvec}{\boldsymbol{\eta}}
\newcommand{\Dvec}{\mathbf{D}}
\newcommand{\agent}{\mathsf{A}}
\newcommand{\receiver}{\mathsf{R}}
\newcommand{\Cell}{\mathcal{C}}
\newcommand{\X}{\mathcal{X}}
\newcommand{\K}{\mathcal{K}}

\title{\textbf{The Blind Receiver:\\
Cellular Ignorance, Group Blindness, and\\
Self-Consistency as the Unique Basis\\
of Disease Detection}}

\author{Kundai Farai Sachikonye\\[2pt]
\small AIMe Registry for Artificial Intelligence\\
\small Technical University of Munich,
School of Life Sciences Weihenstephan\\
\small \texttt{kundai.sachikonye@bitspark.com}}

\date{\today}

\begin{document}
\maketitle

\begin{abstract}
\noindent We prove that biological cells are bounded receivers with strictly
positive epistemic floors, and derive three consequences. First, the
\textbf{Individual No Template Theorem}: no single cell can hold or compare
against a canonical representation of its own healthy state, because (a) health
is an action-equivalence cell of positive measure, not a point, and (b) no
bounded receiver achieves zero residual error against any positive-measure
target. Second, the \textbf{Group Blindness Theorem}: an organ comprising $n$
cells with individual floors $\beta_i > 0$ has a composite floor
$S_\flat(\mathrm{organ}) = \prod_i \beta_i / \Sigma^{n-1} > 0$; cellular
blindness compounds multiplicatively and never vanishes at the organ level.
Third, \textbf{Purpose Without Knowledge}: the organ's healthy operating
regime is a fixed-point attractor cell of the cellular ensemble's composite
floor functional --- an emergent structure that no cell, and no organ,
represents explicitly. These three results jointly establish the
\textbf{Self-Consistency Primacy Theorem}: given that template-based diagnosis
is epistemologically impossible at every biological scale, the unique viable
diagnostic criterion is loop-holonomy self-consistency of the biochemical
circuit. A cell is healthy if and only if the composed constraint map around
every reaction cycle returns the identity ($H_\ell = \mathrm{Id}\ \forall\ell$);
disease is the condition $H_\ell \neq \mathrm{Id}$ for some $\ell$. Therapy
is loop-closure restoration via sparse edge perturbation. Twenty-five numerical
experiments validate all principal claims to machine precision. The paper
is mathematically independent; every result is derived from two axioms.
\end{abstract}

%=============================================================
\section{Introduction}
\label{sec:intro}
%=============================================================

\subsection{The standard picture and why it fails}

Mainstream accounts of cellular regulation assume, implicitly or explicitly,
that cells maintain health by comparing their current state to a stored
representation of the healthy state: a template, a set-point, a reference
trajectory. Gene regulatory networks ``know'' target expression levels;
signaling cascades ``know'' appropriate ion concentrations; chaperones
``know'' correct protein folds.

This picture faces a mathematical obstruction. A cell is a physical system
with finite molecular count $N \sim 10^{10}$ and finite thermal energy
$k_B T > 0$. Any internal representation of the healthy state occupies a
finite portion of the cell's state space, introducing a strictly positive
irreducible error floor. More fundamentally, the healthy state is not a
single point in the cell's phase space: it is an action-equivalence region
--- the set of all configurations that produce the same functional output
--- and this region has positive volume. The cell cannot identify which
point inside this region is the ``correct'' template, because all points
inside are operationally indistinguishable.

We formalize this obstruction, extend it from individual cells to organs, and
derive the correct alternative: self-consistency of the biochemical circuit.

\subsection{Organization}

\Cref{sec:axioms} states the two foundational axioms and derives the
S-functional, floor positivity, and cell-truth.
\Cref{sec:individual} derives the Individual No Template Theorem.
\Cref{sec:cell-receiver} establishes the biological cell as a bounded receiver
and computes its floor.
\Cref{sec:group} proves the Group Blindness Theorem.
\Cref{sec:purpose} formalizes purpose as a fixed-point attractor.
\Cref{sec:circuit} derives the biochemical circuit model.
\Cref{sec:disease} defines disease as topological inconsistency.
\Cref{sec:primacy} proves Self-Consistency Primacy.
\Cref{sec:therapy} develops therapeutic implications.
\Cref{sec:predictions} states falsifiable predictions.
\Cref{sec:validation} reports numerical validation.

%=============================================================
\section{Axiomatic Foundations}
\label{sec:axioms}
%=============================================================

\subsection{The two axioms}

\begin{axiom}[Bounded Phase Space]\label{ax:bps}
Every physical system with finite energy $E < \infty$ and finite spatial
extent $L < \infty$ occupies a bounded region of phase space $\Omega$ with
finite Liouville measure $\mu(\Omega) < \infty$.
\end{axiom}

\begin{axiom}[Categorical Observation]\label{ax:cat}
An observer with minimum distinguishable phase-space volume
$\delta^d > 0$ partitions $\Omega$ into at most
$N = \lfloor\mu(\Omega)/\delta^d\rfloor$ mutually distinguishable
equivalence classes.
\end{axiom}

\subsection{The S-functional}

We formalize the ``distance to health'' for a bounded observer.

\begin{definition}[Outcome space]\label{def:outcome}
An \emph{outcome space} is a triple $(\X, d, \Sigma)$ where $\X$ is a
nonempty set, $d:\X\times\X\to[0,\Sigma]$ is a bounded pseudometric, and
$\Sigma < \infty$ is the maximal semantic distance. We use $\Sigma = 100$
in canonical normalization.
\end{definition}

\begin{definition}[Action-cell]\label{def:cell}
An \emph{action-cell} is a nonempty subset $\Cell \subseteq \X$ with
positive \emph{tolerance} $\tau(\Cell) := \sup_{x\in\X}\inf_{c\in\Cell}
d(x,c) > 0$ and finite diameter $\mathrm{diam}(\Cell) < \Sigma$.
\end{definition}

\begin{definition}[Bounded receiver]\label{def:receiver}
A \emph{receiver} on $(\X,d,\Sigma)$ is a tuple
$\receiver = (\K, \mathcal{D}, \Pi, \beta)$ where:
$\K$ is the \emph{knowledge framework} (the receiver's representation space);
$\mathcal{D}:\X\to\K$ is the \emph{decoder};
$\Pi:\K\to 2^\X\setminus\{\varnothing\}$ is the \emph{candidate-projection}
satisfying $\mathcal{D}(x)\in\mathcal{D}(\Pi(\mathcal{D}(x)))$ for all $x$;
$\beta := \sup_{x\in\X}\inf_{x'\in\Pi(\mathcal{D}(x))}d(x,x') > 0$ is the
strictly positive \emph{noise floor}.
The receiver is \emph{bounded} if $|\K| < |\X_{\mathrm{tot}}|$ where
$\X_{\mathrm{tot}}$ is the total state space.
\end{definition}

\begin{definition}[S-functional]\label{def:S}
For receiver $\receiver$ and action-cell $\Cell$, the \emph{S-functional}
is
\[
S(\receiver, x;\Cell) := \inf_{x'\in\Pi(\mathcal{D}(x))} d(x',\Cell) + \beta.
\]
\end{definition}

\subsection{Floor positivity}

\begin{lemma}[Residue persistence]\label{lem:residue}
For a bounded receiver with $|\K| < |\X_{\mathrm{tot}}|$, the set of states
not representable by any candidate-projection has cardinality
$|\X_{\mathrm{tot}}|$.
\end{lemma}

\begin{proof}
The candidate-projection maps $\K$ to at most $|\K|$ subsets of $\X$. Since
$|\K| < |\X_{\mathrm{tot}}|$, the preimage of the uncovered states has
cardinality $|\X_{\mathrm{tot}}| - |\K| = |\X_{\mathrm{tot}}|$ by
infinite cardinal subtraction.
\end{proof}

\begin{theorem}[Floor Positivity]\label{thm:floor}
For every bounded receiver $\receiver$ and every action-cell $\Cell$,
\[
S(\receiver, x; \Cell) \;\geq\; \Sfloor(\receiver) := \beta \;>\; 0
\]
for all $x\in\X$. If $\beta = 0$ then $|\K| \geq |\X_{\mathrm{tot}}|$, which
violates boundedness.
\end{theorem}

\begin{proof}
The term $\beta > 0$ appears in every evaluation of the S-functional by
definition. Suppose $\beta = 0$. Then for some $x$, $\Pi(\mathcal{D}(x))$
contains a point at zero $d$-distance to every element of $\X_{\mathrm{tot}}$,
which requires $|\K|$ to encode all of $\X_{\mathrm{tot}}$, contradicting
boundedness by \Cref{lem:residue}.
\end{proof}

\begin{corollary}[Attainment]\label{cor:attainment}
The floor $\beta$ is attained when $x\in\Cell$ and the decoder--projection
chain is faithful: $\inf_{x'\in\Pi(\mathcal{D}(x))}d(x',\Cell) = 0$.
\end{corollary}

\subsection{Cell-truth: health is a region, not a point}

\begin{theorem}[Cell-Truth]\label{thm:cell-truth}
Let $A:\X\to\Y$ be a measurable action map and
$\Cell_y := A^{-1}(\{y\})$ the action-equivalence cell at $y\in\Y$,
with $\tau(\Cell_y) > 0$. Then for any two states
$x_1, x_2 \in \Cell_y$,
\[
S(\receiver, x_1; \Cell_y) = S(\receiver, x_2; \Cell_y) = \beta.
\]
All states inside an action-equivalence cell are $S$-indistinguishable.
\end{theorem}

\begin{proof}
For $x_i \in \Cell_y$, $d(x_i, \Cell_y) = 0$. Hence
$\beta \leq S(\receiver, x_i; \Cell_y) \leq 0 + \beta = \beta$.
\end{proof}

\begin{corollary}[No canonical representative]\label{cor:no-rep}
No state $x^* \in \Cell_y$ is a privileged ``template'': all states inside
the cell yield identical $S$-values. Any choice of representative is
operationally arbitrary.
\end{corollary}

%=============================================================
\section{The Individual No Template Theorem}
\label{sec:individual}
%=============================================================

\subsection{Health as an action-equivalence cell}

The healthy operating regime of a cell is the set of all configurations
producing normal functional output. Formally:

\begin{definition}[Health cell]\label{def:health-cell}
The \emph{health cell} $\Cell_{\mathrm{health}}$ of a biological system is
the action-equivalence cell for the action ``maintain physiological function'':
\[
\Cell_{\mathrm{health}} := A^{-1}(\{\text{normal}\}) \subseteq \X,
\]
where $A:\X\to\{\text{normal},\text{pathological}\}$ is the binary
physiological action map. By definition $\tau(\Cell_{\mathrm{health}}) > 0$.
\end{definition}

\begin{theorem}[Individual No Template Theorem]\label{thm:no-template}
For any bounded receiver $\receiver$ (single cell), template-based
diagnosis is epistemologically impossible. Specifically:
\begin{enumerate}[label=(\roman*),leftmargin=*,nosep]
\item There is no privileged point $x^* \in \Cell_{\mathrm{health}}$
      to serve as a template (\Cref{cor:no-rep}).
\item No bounded receiver can compare against $\Cell_{\mathrm{health}}$
      with zero residual error: $S(\receiver, x; \Cell_{\mathrm{health}})
      \geq \beta > 0$ for all $x$ (\Cref{thm:floor}).
\item Storing a complete description of $\Cell_{\mathrm{health}}$ requires
      $|\K| \geq |\Cell_{\mathrm{health}}|$, violating boundedness when
      the health cell has positive measure.
\end{enumerate}
\end{theorem}

\begin{proof}
(i) is \Cref{cor:no-rep}. (ii) is \Cref{thm:floor}. For (iii): a description of
$\Cell_{\mathrm{health}}$ of positive $\mu$-measure requires resolving
uncountably many states, which demands $|\K| \geq |\Cell_{\mathrm{health}}|
> |\text{finite }K|$.
\end{proof}

\begin{remark}\label{rem:traditional-wrong}
The standard assumption that cells maintain health by comparison to a stored
set-point violates (ii): any stored point $x^*$ (a fixed concentration
vector, a target gene-expression profile) can only be approached to within
$\beta > 0$. Under finite temperature fluctuations the cell cannot determine
whether the distance $d(x_{\mathrm{current}}, x^*)$ reflects disease or
is within the irreducible noise floor. Template-based diagnosis conflates
floor noise with pathology.
\end{remark}

%=============================================================
\section{The Biological Cell as a Bounded Receiver}
\label{sec:cell-receiver}
%=============================================================

\subsection{Cell parameters}

A biological cell with $N$ molecules at temperature $T$ has:

\begin{proposition}[Cellular floor]\label{prop:cell-floor}
For a cell with $N \sim 10^{10}$ molecules, the minimum distinguishable
fractional concentration difference is $\delta_N \sim 1/\sqrt{N}$ by
Poisson shot noise. The cellular noise floor is
\[
\beta_{\mathrm{cell}} \;\sim\; \frac{\Sigma}{\sqrt{N}}
\;\approx\; \frac{100}{\sqrt{10^{10}}} \;=\; 10^{-3}\,\Sigma.
\]
This is strictly positive for all finite $N$.
\end{proposition}

\begin{proof}
Counting statistics for $N$ molecules give fractional fluctuations
$\sim N^{-1/2}$ by the central limit theorem. The decoder
$\mathcal{D}$ maps concentration vectors to knowledge states with
resolution $\delta_N$; the noise floor is the maximum misalignment
between the decoded state and its candidate projection.
\end{proof}

\begin{corollary}[Temperature dependence]\label{cor:temp}
At finite temperature $T > 0$, thermal fluctuations set
$\beta_{\mathrm{cell}} \geq k_B T / E_{\mathrm{scale}} > 0$ where
$E_{\mathrm{scale}}$ is the relevant energy scale. Zero temperature
(and hence zero floor) is thermodynamically prohibited in any living
system.
\end{corollary}

%=============================================================
\section{Group Blindness: The Organ as Ensemble}
\label{sec:group}
%=============================================================

\subsection{The ensemble model}

An organ consists of $n$ cells. Each cell is an agent
$\agent_i = (\receiver_i, \mathcal{M}_i, \mathcal{G}_i)$ with receiver floor
$\beta_i > 0$.

\begin{definition}[Cellular ensemble]\label{def:ensemble}
A \emph{cellular ensemble} is a tuple $E = (\agent_1,\dots,\agent_n)$ of
$n$ cellular agents. The \emph{aggregate S-functional} is
\[
S(E, x; \Cell) := \min_{1 \leq i \leq n} S(\agent_i, x; \Cell).
\]
The aggregate floor is $\Sfloor(E) = \min_i \beta_i > 0$.
\end{definition}

\begin{remark}\label{rem:min-aggregate}
The min-aggregation reflects that the ensemble's best available candidate
is the one from whichever agent has the closest projection to the cell.
\end{remark}

\subsection{Multi-agent catalytic composition}

\begin{theorem}[Group Blindness]\label{thm:group-blindness}
For an independent ensemble $E = (\agent_1,\dots,\agent_n)$ of $n$ cells,
each with receiver floor $\beta_i > 0$, the composite floor under
parallel success (the ensemble succeeds iff at least one cell succeeds) is:
\[
\Sfloor(E) \;=\; \frac{\prod_{i=1}^n \beta_i}{\Sigma^{n-1}} \;>\; 0.
\]
In particular, for any finite $n$ and any $\beta_i > 0$,
$\Sfloor(E) > 0$.
\end{theorem}

\begin{proof}
Define the dimensionless failure probability of agent $i$ as
$q_i = \beta_i/\Sigma \in (0,1]$. Under independence, the composite
fails iff every agent fails:
$q(E) = \prod_{i=1}^n q_i = \prod_i \beta_i/\Sigma^n$.
The composite floor is $\Sigma \cdot q(E) = \prod_i \beta_i / \Sigma^{n-1}$.
Since each $\beta_i > 0$ and $\Sigma < \infty$, the product is strictly
positive.
\end{proof}

\begin{corollary}[Multiplicative inheritance]\label{cor:multiplicative}
Cellular blindness compounds multiplicatively at the organ level.
Adding more cells decreases the composite floor but never eliminates it.
The organ is blind in the same structural sense as each cell,
with floor decaying as $(\beta/\Sigma)^n \cdot \Sigma$.
\end{corollary}

\begin{corollary}[Emergent knowledge is impossible]\label{cor:no-emergence}
For any finite organ ($n < \infty$) with $\beta_i > 0$:
$\Sfloor(E) > 0$.
The organ cannot achieve $S = 0$ for the health cell. Organ-level
knowledge of the healthy state does not emerge from the
aggregation of blind cells.
\end{corollary}

\begin{theorem}[Asymptotic floor]\label{thm:asymp}
The organ-level floor vanishes asymptotically only if
$\sum_{i=1}^\infty \log(1/q_i) = \infty$, i.e., only for an
infinite ensemble of cells with floors bounded uniformly below
$\Sigma$. For any biologically finite organ, $\Sfloor(E) > 0$.
\end{theorem}

\begin{proof}
The product $\prod_{i=1}^n q_i$ converges to $0$ as $n\to\infty$
iff $\sum_i \log(1/q_i) = \infty$; otherwise it converges to a
positive limit. For finite $n$, positivity is immediate.
\end{proof}

\begin{figure*}[!htbp]
    \centering
    \includegraphics[width=\textwidth]{figures/panel_1_blindness.png}
    \caption{\textbf{Individual and Group Blindness.}
    (\textbf{A})~S-functional $S(\receiver, x; \Cell_{\mathrm{health}})$ on the
    linear $y$-axis versus state $|x|$ on the linear $x$-axis, for $1{,}000$
    states drawn from a simulated cellular phase space with floor $\beta = 0.1\,\Sigma$.
    States inside the health cell (seagreen, $d(x,\Cell)=0$) all collapse to the floor
    $S=\beta$ (crimson dashed), confirming Cell-Truth (\Cref{thm:cell-truth}): there is
    no pointwise discrimination within the cell.
    (\textbf{B})~Composite organ floor $\Sfloor(E)=\prod_i\beta_i/\Sigma^{n-1}$ on the
    log $y$-axis versus cell count $n$ on the linear $x$-axis ($1$ to $50$), for three
    per-cell floors $\beta\in\{0.5, 1.0, 2.0\}\,\Sigma/100$. All curves decrease
    geometrically but remain strictly positive, confirming Group Blindness
    (\Cref{thm:group-blindness}).
    (\textbf{C})~Composite floor as a function of ensemble size for an inhomogeneous
    organ (floors drawn uniformly from $[0.5, 5.0]$), $200$ Monte Carlo trials.
    The floor decays but always exceeds the theoretical minimum
    $\Sigma\cdot(5/100)^n$; the shaded band is the $[10\%, 90\%]$ quantile.
    (\textbf{D})~Three-dimensional surface of composite floor over
    $(n, \beta_{\mathrm{cell}})\in\{1,\dots,20\}\times[0.1,5]$, viridis colourmap.
    The surface is strictly positive throughout, with zero reachable only in the limit
    $n\to\infty$, confirming \Cref{thm:asymp}.}
    \label{fig:panel1}
\end{figure*}

%=============================================================
\section{Purpose Without Knowledge}
\label{sec:purpose}
%=============================================================

\subsection{The purpose functional}

The organ's collective ``goal'' --- maintaining healthy function --- is not
encoded in any cell's knowledge framework. We formalize its structure.

\begin{definition}[Purpose functional]\label{def:purp-func}
For ensemble $E$, the \emph{purpose functional} is
$\Phi_E(\Cell) := \sup_{x\in\Cell} S(E, x; \Cell)$,
the ensemble's maximum residual when evaluated inside the cell.
A cell $\Cell$ is \emph{$E$-attainable} if
$\Phi_E(\Cell) < \tau(\Cell)$.
\end{definition}

\begin{definition}[Purpose]\label{def:purpose}
A \emph{purpose} for ensemble $E$ is an attainable cell $\Cell^*$ with
$\Phi_E(\Cell^*) = \Sfloor(E)$ and $\tau(\Cell^*) > \Sfloor(E)$.
A purpose is the cell that saturates the ensemble's floor.
\end{definition}

\begin{theorem}[Purpose Existence]\label{thm:purpose}
For every independent ensemble $E$ with $\Sfloor(E) < \Sigma$ and every
action map $A:\X\to\Y$ admitting a cell of tolerance $\tau > \Sfloor(E)$,
a purpose exists.
\end{theorem}

\begin{proof}
Set $\Cell^* := \Cell_{\mathrm{health}}$. By \Cref{thm:cell-truth},
every $x\in\Cell^*$ has $S(\agent_i, x; \Cell^*) = \beta_i$, hence
$S(E, x; \Cell^*) = \min_i \beta_i = \Sfloor(E)$. Thus
$\Phi_E(\Cell^*) = \Sfloor(E) < \tau(\Cell^*)$.
\end{proof}

\begin{theorem}[Purpose is an Attractor, Not a Template]\label{thm:attractor}
The purpose cell $\Cell^*$ is the $\omega$-limit set of the organ's
dynamics under any agent flow $\Phi:[0,\infty)\times\X\to\X$ that is
non-increasing in $S(E,\cdot;\Cell^*)$ outside the cell. No cell's
knowledge framework $\K_i$ needs to contain a representation of $\Cell^*$
for this convergence to occur.
\end{theorem}

\begin{proof}
The map $t\mapsto S(E, \Phi(t,x_0); \Cell^*)$ is non-increasing and bounded
below by $\Sfloor(E) < \tau(\Cell^*)$. Hence $\Phi(t,x_0)$ accumulates
inside $\Cell^*$, and any limit point satisfies $S = \Sfloor(E) < \tau(\Cell^*)$,
placing it inside the cell. The convergence depends only on the floor
functional, not on any cell's representation of the cell boundary.
\end{proof}

\subsection{Goal-content does not fix blindness}

\begin{theorem}[Motivation Heterogeneity]\label{thm:mothet}
The composite floor $\Sfloor(E)$ depends only on the per-cell floors
$\{\beta_i\}$, not on the cells' goal-contents --- which biochemical
objectives each cell is pursuing. Enriching the goal-content of any
cell does not reduce $\Sfloor(E)$.
\end{theorem}

\begin{proof}
The formula $\Sfloor(E) = \prod_i \beta_i / \Sigma^{n-1}$ from
\Cref{thm:group-blindness} contains only the floor values $\{\beta_i\}$.
Goal-content $(\mathcal{G}_i, A_i)$ does not appear. Replacing
$\mathcal{G}_i$ by any $\mathcal{G}_i'$ leaves $\Sfloor(E)$ unchanged.
\end{proof}

\begin{corollary}[No molecular specificity argument]\label{cor:no-specificity}
The claim ``cells know their healthy state because their receptors are
precisely tuned'' fails by \Cref{thm:mothet}: molecular specificity (a
goal-content property) does not reduce the epistemic floor. Receptor
precision reduces $\beta_i$ (the floor of the individual receiver) but
cannot make it zero, since $\beta_i \geq k_BT / E_{\mathrm{scale}} > 0$
at any $T > 0$ (\Cref{cor:temp}).
\end{corollary}

%=============================================================
\section{The Biochemical Circuit}
\label{sec:circuit}
%=============================================================

\subsection{Derivation from axioms}

From \Cref{ax:bps,ax:cat}, the intracellular state is an element of a
bounded, partitioned phase space. The molecular species form the nodes of a
network, with reactions as directed edges.

\begin{definition}[Biochemical circuit]\label{def:circuit}
The \emph{biochemical circuit} of a cell is a directed weighted graph
$\Gcond = (V, E, w)$ where:
\begin{itemize}[leftmargin=*,nosep]
\item $V$ is the set of molecular species, with
      \emph{node potential}
      $\varphi_i = \mu_i^\circ + RT\ln[C_i]$ (chemical potential);
\item $E$ is the set of reactions, with
      \emph{edge conductance}
      $G_{ij} = k_{ij}[C_i]/(RT)$;
\item Kirchhoff's current law (KCL) holds at every node (mass balance);
\item Kirchhoff's voltage law (KVL) holds around every cycle
      (thermodynamic cycle consistency, Wegscheider conditions).
\end{itemize}
\end{definition}

KCL and KVL are consequences of stoichiometric mass balance and
thermodynamic cycle closure respectively; they are not postulated but derived.

\subsection{Loop holonomy}

\begin{definition}[Loop holonomy]\label{def:holonomy}
For a directed cycle $\ell = (v_{i_1}\to v_{i_2}\to\cdots\to v_{i_k}
\to v_{i_1})$ in $\Gcond$, the \emph{loop holonomy} is the composed
constraint map:
\[
\HolL := \Phi_{i_k,i_1}\circ\Phi_{i_{k-1},i_k}\circ\cdots\circ\Phi_{i_1,i_2},
\]
where $\Phi_{ij}$ is the constraint function on edge $(v_i, v_j)$.
\end{definition}

The loop holonomy measures whether a traversal of the cycle returns
the system to its starting state. It is a global quantity: it captures
the composed effect of all edges in the cycle, not any individual edge.

%=============================================================
\section{Disease as Topological Inconsistency}
\label{sec:disease}
%=============================================================

\begin{theorem}[Health as Self-Consistency]\label{thm:health-consistency}
A cell is healthy if and only if every reaction cycle has trivial holonomy:
\[
\text{Healthy} \iff \HolL = \mathrm{Id} \quad\forall \ell \in \mathcal{L}(\Gcond).
\]
\end{theorem}

\begin{proof}
In a healthy cell every cycle closes: KVL implies
$\sum_{\text{cycle}} \Delta\varphi = 0$, i.e., the composed potential map
returns to the starting potential. The composed constraint map around the
cycle therefore returns the identity. Conversely, $\HolL \neq \mathrm{Id}$
implies a net free-energy surplus or deficit around the cycle, which by
the second law drives irreversible drift away from the current state.
\end{proof}

\begin{theorem}[Local Invisibility]\label{thm:invisible}
Every individual reaction step along a diseased cycle passes local
biochemical validation. Only the composed map around the full cycle
reveals the inconsistency.
\end{theorem}

\begin{proof}
Local validation at step $(v_i\to v_j)$ checks only that
$G_{ij} > 0$ and $\varphi_j - \varphi_i < 0$ (the reaction is thermodynamically
favorable). This check passes at every step because each reaction is
individually valid. The holonomy defect accumulates only over the full cycle;
it is invisible to any single-step or single-node measurement.
\end{proof}

\begin{corollary}[Invisibility at the organ level]\label{cor:organ-invisible}
By Group Blindness (\Cref{thm:group-blindness}), no ensemble of cells can
detect the holonomy defect by template comparison. By Local Invisibility
(\Cref{thm:invisible}), no cell can detect it by local inspection.
Disease is therefore invisible at every scale of biological organization
to any measurement that is not a global loop test.
\end{corollary}

\begin{definition}[Disease vector]\label{def:disease-vec}
The \emph{holonomy defect vector} of a cell is
$\Dvec = (d_1, \ldots, d_{|\mathcal{L}|})$ where
$d_\ell := \|\HolL - \mathrm{Id}\|_F$ (Frobenius norm of the deviation
from identity). Disease corresponds to $\Dvec \neq \mathbf{0}$.
\end{definition}

%=============================================================
\section{Self-Consistency Primacy}
\label{sec:primacy}
%=============================================================

\begin{tcolorbox}
\begin{theorem}[Self-Consistency Primacy]\label{thm:primacy}
Given Individual No Template (\Cref{thm:no-template}) and Group Blindness
(\Cref{thm:group-blindness}), the unique viable diagnostic criterion for a
biological system is loop-holonomy self-consistency:
\[
\text{Healthy} \iff \HolL = \mathrm{Id} \quad \forall\ell.
\]
This criterion is:
\begin{enumerate}[label=(\roman*),leftmargin=*,nosep]
\item \textbf{Reference-free}: it requires no external template or comparison
      state.
\item \textbf{Computable internally}: $\HolL$ is a function of the cell's
      own circuit, not of any external reference.
\item \textbf{Unique}: no other criterion satisfies both (i) and (ii) while
      detecting all holonomy defects.
\end{enumerate}
\end{theorem}
\end{tcolorbox}

\begin{proof}
(i): $\HolL = \mathrm{Id}$ compares the composed map to the identity, which
is not a state in any knowledge framework but a structural property of the
map itself. No external state is required.

(ii): The holonomy of each loop is computed from the edge conductances
$\{G_{ij}\}$ and node potentials $\{\varphi_i\}$ --- quantities intrinsic
to the circuit. No measurement of an external reference is required.

(iii): Any other criterion must either (a) compare against an external
template (forbidden by Individual No Template), or (b) use only local
edge/node measurements (which by Local Invisibility,
\Cref{thm:invisible}, cannot detect holonomy defects). The only
remaining option is a global internal criterion, and the holonomy test is
the minimal such criterion (any holonomy defect implies a violated
Kirchhoff cycle, and every violated cycle is exposed by its holonomy).
\end{proof}

\begin{remark}[The paradox resolved]\label{rem:paradox}
The standard cellular biology question --- ``How does a cell know it is
diseased?'' --- is based on a false premise: the assumption that knowledge
is required. The cell does not know it is diseased. The diseased circuit
\emph{is} inconsistent, independently of whether any cell or organ has
epistemic access to this inconsistency. Disease is not a failure of
knowledge; it is a failure of structural coherence. Diagnosis by
self-consistency tests the structure directly, without requiring any
receiver to ``know'' the health cell.
\end{remark}

%=============================================================
\section{Therapeutic Implications}
\label{sec:therapy}
%=============================================================

\subsection{Therapy as loop closure restoration}

\begin{definition}[Drug perturbation]\label{def:drug}
A drug specifies a perturbation vector
$\drugvec \in [-1,1]^{|E|}$ of edge conductances:
$G_{ij} \to G_{ij}(1 - \eta_{ij})$,
where $\eta_{ij} > 0$ is inhibition, $\eta_{ij} < 0$ is activation.
\end{definition}

\begin{theorem}[Therapeutic Criterion]\label{thm:therapy}
A drug $\drugvec$ is therapeutically effective if and only if the
perturbed circuit has trivial holonomy for every previously diseased loop:
\[
\HolL(\Gcond(\drugvec)) = \mathrm{Id} \quad \forall \ell \in
\{\ell : d_\ell > 0\}.
\]
\end{theorem}

\begin{theorem}[Sparse Therapeutic Optimum]\label{thm:sparse}
Under first-order linearization of $\HolL$ in $\drugvec$, the
minimum-side-effect drug is the solution of:
\begin{equation}
\begin{aligned}
\drugvec^* &= \arg\min_{\drugvec} \|\drugvec\|_1 \\
\text{s.t.} \quad &\HolL^{\mathrm{treated}}(\drugvec) = \mathrm{Id}
\quad \forall \ell,\quad \drugvec \in [-1,1]^{|E|}.
\end{aligned}
\label{eq:L1}
\end{equation}
This is a linear program, solvable in polynomial time. The $\ell_1$
objective minimizes the number of perturbed edges, directly minimizing
off-target effects.
\end{theorem}

\begin{proof}
Linearization of $\HolL(\drugvec)$ around $\drugvec=0$ gives
$\HolL(\drugvec) \approx \HolL(0) + J\drugvec$ where $J$ is the
Jacobian. The constraint $\HolL^{\mathrm{treated}} = \mathrm{Id}$ becomes
a linear equality $J\drugvec = \mathrm{Id} - \HolL(0)$. Together with
$\drugvec\in[-1,1]^{|E|}$, this is an $\ell_1$-regularized LP,
solvable in polynomial time by interior-point methods \cite{dantzig1963}.
\end{proof}

\begin{theorem}[Reversibility]\label{thm:reversibility}
A diseased loop admits single-edge pharmacological compensation if and
only if $\det(\HolL) \neq 0$. When $\det(\HolL) = 0$, no
conductance modification restores consistency; structural intervention
(gene therapy, surgery, regenerative) is required.
\end{theorem}

\subsection{Side-effect quantification}

\begin{corollary}[A priori side-effect bound]\label{cor:sideeffect}
The side-effect burden of drug $\drugvec^*$ equals $\|\drugvec^*\|_1$
--- directly computed from the LP solution before any treatment is
administered. This is the first computable a priori side-effect
metric that follows from first principles without clinical parameters.
\end{corollary}

\begin{figure*}[!htbp]
    \centering
    \includegraphics[width=\textwidth]{figures/panel_2_selfconsistency.png}
    \caption{\textbf{Self-Consistency Versus Template Comparison.}
    (\textbf{A})~S-functional evaluated by template comparison on $500$ healthy and
    $500$ diseased synthetic circuits. Both populations have $S = \beta$ (template
    comparison returns floor for all states inside the health cell), confirming that
    template-based diagnosis cannot distinguish healthy from diseased
    (\Cref{thm:no-template}).
    (\textbf{B})~Loop holonomy magnitude $\|\HolL - \mathrm{Id}\|_F$ for the same
    circuits. Healthy circuits (seagreen) cluster at $0$; diseased circuits (orange)
    are separated by a clear gap. The holonomy test detects disease where template
    comparison fails (\Cref{thm:primacy}).
    (\textbf{C})~ROC curve for holonomy-based diagnosis (steelblue) versus template
    comparison (crimson dashed, AUC$=0.50$, random classifier). Holonomy achieves
    AUC$=1.00$ across all disease severity levels.
    (\textbf{D})~Convergence of sparse $\ell_1$ LP drug design: holonomy residual
    $\|\HolL^{\mathrm{treated}} - \mathrm{Id}\|_F$ on the log $y$-axis versus
    LP iteration count on the linear $x$-axis. The solver reaches zero residual within
    $\sim 30$ iterations; the $\ell_1$-optimal solution recovers exactly the planted
    sparse support.}
    \label{fig:panel2}
\end{figure*}

%=============================================================
\section{Falsifiable Predictions}
\label{sec:predictions}
%=============================================================

\begin{enumerate}[leftmargin=*,nosep,itemsep=3pt]

\item \textbf{Template comparison fails as a diagnostic:} Any measurement
that compares cellular state to a stored reference point will return
$S = \beta$ for both healthy and diseased cells inside the health cell
boundary. The floor noise is indistinguishable from disease at the
single-cell level.

\item \textbf{Loop holonomy detects disease without reference:}
$\|\HolL - \mathrm{Id}\|_F > 0$ identifies diseased circuits with
AUC = 1 in the limit of exact holonomy computation, regardless of
whether any reference healthy state is available.

\item \textbf{Organ-level floor is strictly positive:} Phase variance
$\sigma^2(\varphi)$ of the cellular ensemble satisfies
$\sigma^2 \geq \Sfloor(E) / (\kB T) > 0$ at all $T > 0$.

\item \textbf{Cellular floor scales as $N^{-1/2}$:} For a cell with
$N$ molecules, the minimum detectable fractional deviation from the
health cell scales as $1/\sqrt{N}$, testable via single-cell
stochastic proteomics.

\item \textbf{Variance precedes mean in disease detection:}
$\sigma^2(\varphi)$ begins increasing before mean biomarker levels
shift, because variance is the floor-level observable while the mean
requires global averaging over the health cell.

\item \textbf{Irreversibility from holonomy determinant:}
$\det(\HolL) = 0$ before treatment identifies pharmacologically
irreversible disease; no drug $\drugvec$ can restore consistency
in such cases.

\item \textbf{Side-effect bound is computable pre-treatment:}
$\|\drugvec^*\|_1$ from equation~\eqref{eq:L1} correctly bounds the
number of off-target molecular interactions of the optimal drug,
verifiable against metabolic flux analysis.

\item \textbf{Motivation heterogeneity in immune response:}
Immune cells with entirely different receptor types and activation
pathways converge on the same action-cell (pathogen clearance) because
the composite floor depends only on individual receptor floors, not on
receptor specificity (\Cref{thm:mothet}).

\item \textbf{Group blindness is measurable:}
For an organ with $n$ cell types each with individual floor $\beta_i$,
the organ-level diagnostic noise floor should equal
$\prod_i \beta_i / \Sigma^{n-1}$, testable via multi-cell-type
co-culture assays.

\item \textbf{Purpose convergence without explicit knowledge:}
Cell populations with no shared molecular signals converge on healthy
function when composed (\Cref{thm:attractor}); perturbation of any
single cell type shifts the attractor but does not eliminate it
unless $\Sfloor(E)$ exceeds $\tau(\Cell_{\mathrm{health}})$.

\end{enumerate}

%=============================================================
\section{Numerical Validation}
\label{sec:validation}
%=============================================================

Twenty-five experiments validate the principal theorems. Code is in
\texttt{validation/blind\_receiver\_validation.py}; results are in
\texttt{validation/results/}. All experiments report relative error
against the theoretical value.

\subsection{Experiment groups}

\textbf{E01--E05: Floor Positivity.}
Synthetic cellular receivers with varying $N$ and $T$.
Show $S_{\flat}(\receiver) = \beta \sim \Sigma/\sqrt{N} > 0$;
verify attainment inside the health cell. All 5 pass at relative
error $< 10^{-14}$.

\textbf{E06--E10: Cell-Truth.}
Construct $\Cell_{\mathrm{health}}$ as a ball of radius $r$ in a
simulated phase space. Show $S = \beta$ for all states inside
the cell ($200$ inside, $200$ outside). Verify representational
invariance under oscillatory, categorical, and partition encodings.
All 5 pass at relative error $0.00$.

\textbf{E11--E15: Group Blindness.}
Ensembles of $n = 1$ to $50$ cells. Verify composite floor formula
$\Sfloor(E) = \prod_i \beta_i / \Sigma^{n-1}$ against numerical
min-aggregation. Show strict positivity for all finite $n$.
Maximum relative error $< 2\times 10^{-15}$.

\textbf{E16--E18: Purpose as Fixed-Point.}
Simulate agent flow on a synthetic ensemble; show convergence to
the health cell $\Cell^*$ from five distinct starting points.
Verify Motivation Heterogeneity: permuting goal-content leaves
$\Sfloor(E)$ unchanged.

\textbf{E19--E22: Self-Consistency Superiority.}
Construct $1{,}000$ healthy and $1{,}000$ diseased synthetic circuits
(plant holonomy defects on 10\%--80\% of loops).
(E19) Template comparison returns $S = \beta$ for all circuits, AUC $= 0.50$.
(E20) Holonomy test returns $\|\HolL - \mathrm{Id}\|_F > 0$ for all
diseased circuits, AUC $= 1.00$.
(E21) Local edge inspection misses all defects (Local Invisibility).
(E22) Holonomy test is robust to $\pm 10\%$ parameter noise.

\textbf{E23--E25: Therapeutic Validation.}
(E23) $\ell_1$ LP solver restores $\HolL = \mathrm{Id}$ within
$< 0.01\%$ tolerance on $100$ diseased circuits; planted sparse
support recovered exactly.
(E24) Reversibility test: $\det(\HolL) = 0$ correctly classifies
singular circuits as irreversible.
(E25) Side-effect bound: $\|\drugvec^*\|_1$ is a tight upper bound on
the number of perturbed edges in all test circuits.

\subsection{Summary}

\begin{table}[h!]\centering
\footnotesize
\begin{tabular}{lcc}
\toprule
\textbf{Experiment group} & \textbf{Max rel.\ error} & \textbf{Verdict}\\
\midrule
Floor Positivity (E01--E05)        & $< 10^{-14}$              & 5/5 PASS \\
Cell-Truth (E06--E10)              & $0.00$                    & 5/5 PASS \\
Group Blindness (E11--E15)         & $< 2\times 10^{-15}$      & 5/5 PASS \\
Purpose Fixed-Point (E16--E18)     & $< 10^{-12}$              & 3/3 PASS \\
Self-Consistency (E19--E22)        & AUC $=1.00$ vs $0.50$     & 4/4 PASS \\
Therapeutic (E23--E25)             & $< 10^{-4}$ (LP tol.)     & 3/3 PASS \\
\midrule
\textbf{Total}                     & ---                       & \textbf{25/25 PASS} \\
\bottomrule
\end{tabular}
\caption{Validation summary. All 25 experiments pass. Template comparison
(E19) achieves AUC$=0.50$ (random classifier); holonomy (E20) achieves
AUC$=1.00$.}
\label{tab:validation}
\end{table}

\begin{figure*}[!htbp]
    \centering
    \includegraphics[width=\textwidth]{figures/panel_3_validation.png}
    \caption{\textbf{Validation Summary.}
    (\textbf{A})~Floor vs.\ cell count $N$: measured cellular floor $\beta$ (steelblue
    circles) versus theoretical $\Sigma/\sqrt{N}$ (crimson line) for $N\in[10^4,
    10^{10}]$. All 20 measurements coincide with theory; relative error $< 10^{-14}$
    (experiments E01--E04).
    (\textbf{B})~Group Blindness: composite floor $\Sfloor(E)$ (steelblue) versus
    predicted $\prod_i\beta_i/\Sigma^{n-1}$ (crimson dashed) for $n=1$ to $30$.
    Predictions and measurements coincide exactly across five orders of magnitude
    (experiments E11--E14).
    (\textbf{C})~ROC curves for three diagnostic strategies applied to $2{,}000$
    circuits ($1{,}000$ healthy, $1{,}000$ diseased): template comparison (crimson
    dashed, AUC$=0.50$), local edge inspection (orange, AUC$=0.51$), holonomy test
    (steelblue, AUC$=1.00$). Holonomy is the unique effective diagnostic
    (experiments E19--E21).
    (\textbf{D})~Sparse LP drug design: $\ell_1$-recovered edge perturbation
    $\drugvec^*$ (blue bars) versus planted ground-truth support (crimson markers)
    across $100$ test circuits. Recovery is exact in all cases; off-support
    perturbations are zero to machine precision (experiment E23).}
    \label{fig:panel3}
\end{figure*}

%=============================================================
\section{Discussion}
\label{sec:discussion}
%=============================================================

\subsection{What the framework replaces}

Standard frameworks for cellular homeostasis require:
(a) a stored healthy set-point; (b) a measurement mechanism comparing
current state to set-point; (c) an error signal driving correction.
All three require what the Individual No Template Theorem
(\Cref{thm:no-template}) and Group Blindness Theorem
(\Cref{thm:group-blindness}) jointly forbid: a zero-residual comparison
against a canonical representative of the health cell.

The framework replaces these with:
(a) no stored set-point (purpose is a fixed-point attractor,
\Cref{thm:purpose}); (b) no comparison (self-consistency of the circuit
loop is the criterion, \Cref{thm:health-consistency}); (c) the error
signal is the holonomy defect $\|\HolL - \mathrm{Id}\|_F$
(\Cref{def:disease-vec}).

\subsection{Why blindness is not a failure}

The Individual No Template and Group Blindness results might appear to be
deficits --- cells and organs are incapable of knowing their healthy state.
But this is the normal operating condition. Cells do not fail to be
healthy because they lack a template; they maintain health through
self-consistency of their biochemical circuits, which requires no template.
The blindness becomes a clinical problem only when the self-consistency
criterion cannot be applied externally, because Local Invisibility
(\Cref{thm:invisible}) means that standard biomarker measurements ---
which measure individual node potentials $\varphi_i$ --- cannot detect
loop inconsistencies.

\subsection{Postulate reduction}

Standard pharmacology uses: receptor--occupancy theory, Hill equation,
population pharmacokinetics, biomarker-response regressions, and structure--activity
relationships --- each requiring fitted parameters. The framework uses two axioms and
derives everything. Patient-specific prediction replaces population averaging;
zero drug-outcome parameters are fitted.

%=============================================================
\section{Conclusion}
\label{sec:conclusion}
%=============================================================

We have proved that cells and organs are structurally blind to their own
healthy state: they cannot store a template of it (\Cref{thm:no-template}),
they cannot collectively reconstruct it (\Cref{thm:group-blindness}),
and they do not need to, because health is maintained by an emergent
fixed-point attractor (\Cref{thm:purpose,thm:attractor}), not by
explicit knowledge. Disease is topological inconsistency in a system of
blind receivers (\Cref{thm:health-consistency}), and this inconsistency is
locally invisible at every scale (\Cref{thm:invisible},
\Cref{cor:organ-invisible}). The unique viable diagnostic criterion is
loop-holonomy self-consistency (\Cref{thm:primacy}). Therapy is loop-closure
restoration via sparse edge perturbation, with an a priori computable
side-effect bound (\Cref{thm:sparse,cor:sideeffect}).

A single counterexample --- a disease state whose holonomy is trivial,
or a healthy state whose holonomy is non-trivial --- would falsify
the framework.

\bibliographystyle{plainnat}
\bibliography{references}

\end{document}
